\chapter{Structural Monitoring}
\label{ch:assignment7}

\begingroup\itshape Assignment 7 --- 8 points\endgroup

\section{Inclinometer array for warping identification}

Conventional inclinometers are rotation sensors used to measure the pitching and rolling of a vessel. High-end electronic inclinometers can provide resolutions in the order of microradians, hence may also be used to measure the twist angle at any point along the ship. If an array of inclinometers is used as a monitoring system, how can the contributions of free and constrained warping be identified with the help of the generalised torsion equation (inertia effects are neglected)? Furthermore, how can the location of the maximum warping stress be identified? \textbf{(4 pts)}

\bigskip
An array of $n$ inclinometers mounted along the length (e.g.\ on the centre line at every bulkhead and at mid-hold) measures the roll angle $\theta_k(t)$ at stations $x_k$. Subtracting the rigid-body roll (the common part, e.g.\ the value at a reference station or the mean) leaves the elastic twist $\theta(x_k)-\theta(x_{ref})$ along the ship. Fitting a smooth function (piecewise polynomial or the closed-form solution \eqref{eq:gensol} with $C_{1i},C_{2i}$ as unknowns per section) to these samples gives $\theta(x)$ and, by differentiation, the twist rate $\theta'(x)$, the curvature $\theta''(x)$ and $\theta'''(x)$. The generalised torsion equation \eqref{eq:gte} then separates the two load paths directly:
\begin{equation}
  M_t(x)=\underbrace{GI_p\,\theta'(x)}_{\text{free (St.\,Venant) warping}}\;\underbrace{-\,EI_w\,\theta'''(x)}_{\text{constrained warping torque } M_\omega},\qquad
  B_\omega(x)=-EI_w\,\theta''(x)\ \ \text{(bimoment)}.
\end{equation}
With the known section properties ($I_p$, $I_w$ from the drawings, as in Table~\ref{tab:torsion-constants}) the measured $\theta'$ gives the St.\,Venant part and the measured $\theta'''$ the warping part; where $\theta'$ is constant (linear $\theta$, $\theta''=\theta'''=0$) warping is free and the whole torque is St.\,Venant, whereas a $\theta'$ that decreases towards a bulkhead ($\theta''\ne0$) reveals the constrained-warping contribution, in the form of the boundary layers of Figure~\ref{fig:a5BUBUH}. In practice $\theta'''$ from three numerical differentiations is noisy, so the inclinometer data should be fitted with the analytical solution per section (which has only the constants $C_{1i}$, $C_{2i}$ and the unknown $M_t$) rather than differentiated numerically; the fit also identifies the actual torque $M_t(x)$ and, from the size of the boundary layers, the effective restraint provided by the bulkheads. Because inertia effects are neglected, the quasi-static equation can be applied at every time instant.

The warping normal stress is $\sigma_{xx}(s,x)=-E\,\omega(s)\,\theta''(x)$. Its maximum over the length therefore lies where $|\theta''|$ -- the curvature of the twist, i.e.\ the slope of the measured $\theta'(x)$ -- is largest, and over the section at the point of largest $|\omega|$ (the hatch coaming / top of the side shell for an open section, $\omega=\pm13.3a^2$ in Assignment~2.5). From the fitted $\theta'(x)$ the positions of maximum $|\theta''|$ are found: for the constrained holds these are the sections adjacent to the transverse bulkheads (and the engine-room front bulkhead), where $\theta'$ changes fastest; in the middle of a hold $\theta''=0$ and the warping stress vanishes. Thus the inclinometer array locates the hot spots without strain measurements, and the bimoment $B_\omega=-EI_w\theta''$ there gives the stress level $\sigma_{xx,max}=B_\omega\omega_{max}/I_w$, which can be checked against the values predicted in Assignment~5 (e.g.\ $|\theta''|_{max}=1.6\times10^{-4}$~m$^{-2}$ next to the bulkheads for B1).

\bigskip
\noindent\textbf{Final answer.} Fit the elastic twist $\theta(x)$ from the inclinometer array with the generalised torsion solution; $GI_p\theta'$ is the free-warping (St.\,Venant) torque and $-EI_w\theta'''$ the constrained-warping torque. The maximum warping stress occurs where $|\theta''|$ is largest -- at the sections next to the (bulkhead) warping restraints -- and at the point of maximum $|\omega|$ in that section.

\section{Long-base strain gauges}

Consider a cross-section instrumented for longitudinal normal strain monitoring on the inner shell (Figure~12 of the brief). The total strain is the combination of the horizontal bending, vertical bending and warping strain, $\varepsilon_{xx,tot}=\varepsilon_{xx,vb}+\varepsilon_{xx,hb}+\varepsilon_{xx,warp}$. In fatigue assessment, a multiaxial stress state may significantly reduce the lifetime of a structure. Through a monitoring system composed of long-base strain gauges, the longitudinal normal strain at any point along the inner shell of a cross section can be measured. Discuss how the torsion and vertical and horizontal bending responses may be separated using multiple long-base strain gauges in the cross-section. Should these strain gauges be placed on a cross-section close to a bulkhead or away from it? \textbf{(4 pts)}

\bigskip
At a point $(y_k,z_k)$ of the section with sectorial coordinate $\omega_k$ (pole at the shear centre, principal sectorial origin), beam theory gives
\begin{equation}
  \varepsilon_{xx}(y_k,z_k)=\frac{N}{EA}+\frac{M_y}{EI_{yy}}\,z_k-\frac{M_z}{EI_{zz}}\,y_k+\frac{B_\omega}{EI_w}\,\omega_k ,
  \label{eq:strainsep}
\end{equation}
where $z$ is measured from the neutral axis of vertical bending, $y$ from the centre line, and the axial force $N$ is included for generality (it is zero for pure hull-girder loading but not necessarily for temperature effects). Each gauge $k$ yields one linear equation in the unknown generalised loads $(N,M_y,M_z,B_\omega)$, with known coefficients $(1/A,\ z_k/I_{yy},\ -y_k/I_{zz},\ \omega_k/I_w)$ from the section geometry. With at least four gauges at suitably chosen positions (more for a least-squares fit) the system is solved at every time step and the three responses are separated. The choice of positions exploits the symmetries of the U-shaped section (Figure~12): vertical bending strain is \emph{symmetric} in $y$ and varies linearly with $z$; horizontal bending strain is \emph{antisymmetric} in $y$ and independent of $z$; warping strain is \emph{antisymmetric} in $y$ (like $\omega$ in Figure~\ref{fig:omega}) but varies linearly with $z$ along the sides and changes sign at $z=e$, and is also antisymmetric along the bottom. A practical arrangement on the red surfaces of Figure~12 is a pair of gauges high on the port and starboard inner side shells ($\pm y_s$, $z_1$) and a pair low on the sides or at the bottom corners ($\pm y_s$, $z_2$):
\begin{itemize}
  \item sum of a port/starboard pair: $\varepsilon_P+\varepsilon_S=2\left(N/EA+M_yz/EI_{yy}\right)$ -- the antisymmetric (horizontal bending and warping) parts cancel; two such sums at $z_1\ne z_2$ give $M_y$ (and $N$);
  \item difference of a pair: $\varepsilon_P-\varepsilon_S=2\left(-M_zy_s/EI_{zz}+B_\omega\omega(z)/EI_w\right)$ -- the symmetric parts cancel; since $\omega$ varies with $z$ and $y_s$ does not, the differences at two heights $z_1$, $z_2$ separate $M_z$ from $B_\omega$ (the warping part changes sign across $z=e$, the horizontal bending part does not); a gauge near $z=e$, where $\omega\approx0$, measures horizontal bending without warping contamination, and a gauge at the top of the side, where $|\omega|$ is largest, is the most sensitive to warping.
\end{itemize}
The separation requires that the section is instrumented at one $x$-position with all gauges (same $M_y$, $M_z$, $B_\omega$), that local effects (plate bending between stiffeners, hatch-corner concentrations) are averaged out -- which is exactly what long-base gauges do -- and that $\omega_k$, $I_w$ etc.\ are known from the structural drawings. Combining the strain-derived $B_\omega(x)$ with the inclinometer-derived $\theta''(x)$ of Question~7.1 provides a consistency check ($B_\omega=-EI_w\theta''$).

\emph{Close to or away from a bulkhead?} The warping strain is proportional to the bimoment, which is largest next to a warping restraint and zero at mid-hold (Assignment~5): a section close to a bulkhead therefore gives the largest, well-conditioned warping signal and coincides with the fatigue-critical location (maximum $\sigma_{xx,warp}$ superimposed on the bending stresses -- the multiaxial state the question refers to). Directly at the bulkhead, however, the strain field is disturbed by the local stiffness of the bulkhead, its brackets and the hatch-end coaming (three-dimensional stress state, shear lag), and beam theory \eqref{eq:strainsep} is not valid. The recommended position is therefore \emph{close to the bulkhead but outside its local influence zone} -- of the order of half the section depth to one depth away (a few metres), i.e.\ well within the boundary layer $1/\lambda$ for an open hold (which extends over the full hold for BU) but away from the local details. Bending moments vary slowly along the hold, so they are not affected by this choice. For a complete picture of the fatigue loading a second, mid-hold section (where warping is nearly zero) is useful to isolate the bending responses and to calibrate the warping identification.

\bigskip
\noindent\textbf{Final answer.} Use $\ge4$ long-base gauges in one section at symmetric port/starboard positions and two heights; sums of pairs isolate vertical bending (and axial force), differences at two heights separate horizontal bending from warping via the known $\omega(s)$ and $I_w$ (least-squares solution of \eqref{eq:strainsep}). Place the instrumented section close to a bulkhead (largest bimoment, fatigue-critical) but outside the local disturbance zone of the bulkhead, with a mid-hold section as bending reference.
